\subsection{Conservative segmented-rotation comparison}
\label{sec:rotation-followup}

The follow-up uses a clearance of \(0.03125\) model units, a distance
tolerance of \(2\times10^{-4}\), a time target of \(5\times10^{-4}\) on
\([0,1]\), and a cap of 200,000 search tasks. Each configuration is
measured nine times after warm-up with alternating or rotating run order.
An 80-digit analytic evaluation provides the reference for the generated
geometries. The production pairs use the same independent adaptive OCCT
distance cross-check for all three methods, with a \(10^{-6}\) mm numerical
margin and a time resolution of \(2\times10^{-5}\).

For a segment of duration \(h\), the chord between the two exact endpoint
poses differs from the fixed-axis trajectory by at most \(M h^2/8\),
where \(M=\omega^2\max_i\|p_{i,\perp}\|\) bounds control-point acceleration.
The bound extends to each patch point by convexity. The segmented wrapper
splits until the chord bound fits the distance tolerance, then encloses the
rounding errors in endpoint position and velocity. It adds the resulting padding
and the requested clearance to the official Chen OBB projection test,
scaled by each axis norm. The wrapper visits segments in temporal order.
It subdivides retained time intervals before running the official
parameter search and verifies proposed rejections with outward projection
bounds on the chord nets. This check can refine the moving B\'ezier
net. Entering upper bounds require an original-path point pair.
The timing includes segmentation, projection verification, and witness
evaluation.

The fixtures share the same untrimmed biquadratic control nets. A moving
rectangle or curved graph tilts about the \(y\)-axis above a fixed plane.
At first entry the minimum vertical separation is
\(H-\sin(\omega t)\), with the closest moving point at \(x=1\)
and, for the graph, on its minimum-height centerline. One entry
occurs strictly between safe endpoints. The near miss stays above the
threshold by \(0.001\) model units at closest approach.

\begin{table}[!htbp]
\centering\small
\caption{Matched original-trajectory results. Every hit bracket contains
the analytic entry. Width is the largest across nine repeats.}
\label{tab:rotation-followup}
\begin{tabular}{lrrrrr}
\toprule
 & & \multicolumn{2}{c}{Median query (ms)} & \multicolumn{2}{c}{Bracket width}\\
\cmidrule(lr){3-4}\cmidrule(lr){5-6}
Fixture & Reference & Chen chords & Native & Chen chords & Native\\
\midrule
\RotationFollowupRows
\bottomrule
\end{tabular}
\end{table}

Both methods meet the time target on these fixtures. The curve forces
more segment and parameter checks in the chord method than a planar
fixture does. The native search certifies the rotating near miss with its
root directional bound. This comparison covers the padded and audited
rotation wrapper. The zero-distance linear comparison in
Table~\ref{tab:chen-common} measures the upstream implementation directly.

\subsection{Counterbore and a second production STEP source}
\label{sec:counterbore-followup}

The generated counterbore has mouth and throat radii of \(1.375\) mm
and \(0.8125\) mm, with an internal shoulder at their junction and a
pin of radius \(0.375\) mm. Different lateral offsets lead to a clear passage, a first
entry at the shoulder, or a first entry at the mouth. For each circular
rim, the reference pairs the leading pin cap with that rim and combines
the radial and axial gaps. At the two entry offsets the radial gap at the
active rim is \(0.0234375\) mm.

The production replication uses the Voron Trident Front Idler B
assembly \cite{vorontrident}. The extractor keeps original AP214
geometry and occurrence transforms within the idler subtree. We test the
M5 spacer and an F695 bearing against their original M5 screw, and the
M3 washer against its screw. Each approach follows the screw axis with a
transverse displacement of \(0.10\) mm. The three variants share the
same trim handling, search, clearance and reference. They use a swept
box, directional time bounds, or directional time bounds with radial
and axial exclusion.

\begin{table}[!htbp]
\centering\small
\caption{Three-way ablation. The reference column gives the analytic
entry for the counterbore and the lower endpoint of the OCCT bracket for
the Voron interfaces. A star marks a budget-limited result or a bracket
wider than the time target.}
\label{tab:counterbore-followup}
\begin{tabular}{lrrrrr}
\toprule
 & & \multicolumn{3}{c}{Median query (ms)} & \\
\cmidrule(lr){3-5}
Fixture & Reference & Swept box & Directional & With radial/axial & Tasks with bound\\
\midrule
\HoleFollowupRows
\bottomrule
\end{tabular}
\end{table}

The radial and axial filter leaves 23--31 tasks for each counterbore
query, compared with 13,754--44,395 for directional clipping alone.
All three variants agree with the analytic references. For the Voron
washer and bearing, both directional variants meet the time target,
while the swept-box variant reaches the cap with wider witnessed
brackets. The radial and axial filter reduces task counts modestly.
All three methods exhaust their budget on the M5 spacer. The best
certified prefix reaches \(0.523320\), and the independent numerical
reference reaches \(0.975586\) before its own cap with no upper witness.

\subsection{Independent trim-edge certificate audit}
\label{sec:edge-audit}

For a plane with a unit circular trim and a parallel disk of radius
\(0.25\), the disk stays inside the hole in projection. At closest distance
the plate point lies in the interior of the inner trim edge at
\((1,0,0)\). The analytic
distance is \(\sqrt{g^2+(0.375-t)^2}\), with \(g\) the in-plane gap.
Offsets place \(g\) well inside the threshold, \(2^{-26}\) below it,
exactly on it, or \(2^{-23}\) above it. The narrow entry interval is
about \(6.10\times10^{-5}\) wide and both trajectory endpoints are safe
in every case.

The independent verifier compares the returned prefixes and witnessed
times with the analytic reference and reconstructs the residual from
the exported Cartesian coordinate intervals. It checks those intervals
against the two plane equations and the disk and trim radii. The
threshold equality at \(t=0.375\) is an isolated contact. A prefix with no
feasible upper witness leaves its remaining interval unresolved.

\begin{table}[!htbp]
\centering\small
\caption{Independent trim-edge audit. The point-pair residual is an
upper bound recomputed from the archived coordinate intervals.}
\label{tab:edge-audit}
\begin{tabular}{lrrrrr}
\toprule
Case & Reference & Outcome & Safe prefix & Width & Residual\\
\midrule
\WitnessFollowupRows
\bottomrule
\end{tabular}
\end{table}

Both crossing cases return edge witnesses and contain their analytic
entry times. The exact graze remains unresolved, with its prefix ending
\(4.77\times10^{-7}\) before the reference time. The near miss certifies
the entire interval at a minimum margin of \(2^{-23}\) mm.
